\documentclass[12pt]{article}
\usepackage{amsmath, amssymb}
\usepackage{geometry}
\geometry{margin=1in}
\setlength{\parskip}{0.8em}
\setlength{\parindent}{0pt}

\begin{document}

\begin{center}
\Large \textbf{CSE400 – Fundamentals of Probability in Computing} \\
\vspace{0.2cm}
\large \textbf{Lecture 5: Bayes’ Theorem, Random Variables, and Probability Mass Function}
\end{center}

\hrule
\vspace{0.5cm}

\section*{1. Bayes’ Theorem}

\subsection*{1.1 Weighted Average of Conditional Probabilities}

Let $A$ and $B$ be events.

We may express event $A$ as:
\[
A = AB \cup AB^c
\]

This is because, for an outcome to be in $A$, it must either:
\begin{itemize}
    \item be in both $A$ and $B$, or
    \item be in $A$ but not in $B$.
\end{itemize}

Since $AB$ and $AB^c$ are mutually exclusive, by \textbf{Axiom 3}:
\[
\Pr(A) = \Pr(AB) + \Pr(AB^c)
\]

Using conditional probability:
\[
\Pr(A) = \Pr(A|B)\Pr(B) + \Pr(A|B^c)\bigl[1 - \Pr(B)\bigr]
\]

\textbf{Interpretation:}  
The probability of event $A$ is a \textbf{weighted average of conditional probabilities}, where the weights are the probabilities of the events on which we condition.

\subsection*{1.2 Learning by Example – Example 3.1 (Part 1/2)}

\textbf{Problem Statement:}  
An insurance company divides people into two classes:
\begin{itemize}
    \item Accident-prone
    \item Not accident-prone
\end{itemize}

Statistics show:
\begin{itemize}
    \item An accident-prone person has an accident within one year with probability $0.4$.
    \item A non-accident-prone person has an accident within one year with probability $0.2$.
    \item $30\%$ of the population is accident-prone.
\end{itemize}

Find the probability that a new policyholder will have an accident within one year.

\textbf{Solution:}

Let:
\begin{itemize}
    \item $A_1$: policyholder has an accident within one year
    \item $A$: policyholder is accident-prone
\end{itemize}

\[
\Pr(A_1) = \Pr(A_1|A)\Pr(A) + \Pr(A_1|A^c)\Pr(A^c)
\]
\[
= (0.4)(0.3) + (0.2)(0.7) = 0.26
\]

\subsection*{1.3 Learning by Example – Example 3.1 (Part 2/2)}

\textbf{Problem Statement:}  
Suppose a new policyholder \textbf{has} an accident within a year.  
Find the probability that the policyholder is accident-prone.

\textbf{Solution:}
\[
\Pr(A|A_1) = \frac{\Pr(AA_1)}{\Pr(A_1)}
= \frac{\Pr(A)\Pr(A_1|A)}{\Pr(A_1)}
\]
\[
= \frac{(0.3)(0.4)}{0.26} = \frac{6}{13}
\]

\subsection*{1.4 Law of Total Probability}

Suppose:
\begin{itemize}
    \item $B_1, B_2, \ldots, B_n$ are mutually exclusive events
    \item $\bigcup_{i=1}^{n} B_i = B$
\end{itemize}

Exactly one of the events $B_1, B_2, \ldots, B_n$ must occur.

Writing:
\[
A = \bigcup_{i=1}^{n} AB_i
\]

Since $AB_i$ are mutually exclusive:
\[
\Pr(A) = \sum_{i=1}^{n} \Pr(AB_i)
= \sum_{i=1}^{n} \Pr(A|B_i)\Pr(B_i)
\]

This is called the \textbf{Law of Total Probability (Formula 3.4)}.

\subsection*{1.5 Bayes Formula (Proposition 3.1)}

Using:
\[
\Pr(AB_i) = \Pr(B_i|A)\Pr(A)
\]

We obtain:
\[
\Pr(B_i|A) =
\frac{\Pr(A|B_i)\Pr(B_i)}
{\sum_{j=1}^{n} \Pr(A|B_j)\Pr(B_j)}
\]

Where:
\begin{itemize}
    \item $\Pr(B_i)$ is the \textbf{a priori probability}
    \item $\Pr(B_i|A)$ is the \textbf{posteriori probability}
\end{itemize}

\subsection*{1.6 Learning by Example – Example 3.2 (Card Problem)}

Three cards:
\begin{itemize}
    \item Card 1: red--red
    \item Card 2: black--black
    \item Card 3: red--black
\end{itemize}

One card is randomly selected and placed on the ground.  
The upper side is red.

Find the probability that the other side is black.

\textbf{Solution:}

Let:
\begin{itemize}
    \item $RR$: all-red card
    \item $BB$: all-black card
    \item $RB$: red--black card
    \item $R$: upturned side is red
\end{itemize}

\[
\Pr(RB|R) =
\frac{\Pr(R|RB)\Pr(RB)}
{\Pr(R|RR)\Pr(RR) + \Pr(R|RB)\Pr(RB) + \Pr(R|BB)\Pr(BB)}
\]

\[
= \frac{(1/2)(1/3)}{(1)(1/3) + (1/2)(1/3) + (0)(1/3)} = \frac{1}{3}
\]

\section*{2. Random Variables}

\subsection*{2.1 Motivation and Concept}

When an experiment is performed, interest often lies in a \textbf{function of the outcome}, not the outcome itself.

Examples:
\begin{itemize}
    \item In dice tossing, focus is on the sum of dice.
    \item In coin tossing, focus is on the number of heads.
\end{itemize}

\textbf{Definition:}  
A random variable is a real-valued function defined on the sample space.

\[
X : \Omega \rightarrow \mathbb{R}
\]

The lecture restricts attention to \textbf{discrete random variables}, whose values form a finite or countably infinite subset of $\mathbb{R}$.

\subsection*{2.2 Distribution of a Random Variable}

Two key components:
\begin{enumerate}
    \item The set of values the random variable can take
    \item The probabilities with which it takes those values
\end{enumerate}

For any value $a$:
\[
\{\omega \in \Omega : X(\omega) = a\}
\]
is an event.

The probability:
\[
\Pr(X = a)
\]
defines the distribution of the random variable.

\subsection*{2.3 Example – Tossing 3 Fair Coins}

Let $Y$ be the number of heads.

Possible values: $0,1,2,3$

\[
\Pr(Y=0)=\frac{1}{8}, \quad
\Pr(Y=1)=\frac{3}{8}, \quad
\Pr(Y=2)=\frac{3}{8}, \quad
\Pr(Y=3)=\frac{1}{8}
\]

Since $Y$ must take one of these values:
\[
1 = \sum_{i=0}^{3} \Pr(Y=i)
\]

\section*{3. Probability Mass Function (PMF)}

\subsection*{3.1 Concept}

A random variable that can take on at most a countable number of values is called \textbf{discrete}.

Let $X$ be a discrete random variable with range:
\[
R_X = \{x_1, x_2, x_3, \ldots\}
\]

The function:
\[
p(x_k) = \Pr(X = x_k)
\]
is called the \textbf{Probability Mass Function (PMF)} of $X$.

Since $X$ must take one of these values:
\[
\sum_k p(x_k) = 1
\]

\subsection*{3.2 PMF Example}

Given:
\[
p(i) = c \frac{\lambda^i}{i!}, \quad i = 0,1,2,\ldots
\]

Find $\Pr(X=0)$ and $\Pr(X>2)$.

Since:
\[
\sum_{i=0}^{\infty} p(i) = 1
\Rightarrow
c \sum_{i=0}^{\infty} \frac{\lambda^i}{i!} = 1
\]

Using:
\[
e^{\lambda} = \sum_{i=0}^{\infty} \frac{\lambda^i}{i!}
\]

We get:
\[
c e^{\lambda} = 1 \Rightarrow c = e^{-\lambda}
\]

Thus:
\[
\Pr(X=0) = e^{-\lambda}
\]

\[
\Pr(X>2) = 1 - \Pr(X \le 2)
= 1 - \bigl[\Pr(X=0) + \Pr(X=1) + \Pr(X=2)\bigr]
\]

\vspace{0.5cm}


\end{document}
